\documentclass[12pt,a4paper]{article}
\usepackage[utf8]{inputenc}
\usepackage[T1]{fontenc}
\usepackage{geometry}
\usepackage{listings}
\usepackage{xcolor}
\usepackage{hyperref}

\geometry{margin=1in}

\definecolor{codegreen}{rgb}{0,0.6,0}
\definecolor{codegray}{rgb}{0.5,0.5,0.5}
\definecolor{codepurple}{rgb}{0.58,0,0.82}
\definecolor{backcolour}{rgb}{0.95,0.95,0.92}

\lstdefinestyle{mystyle}{
    backgroundcolor=\color{backcolour},   
    commentstyle=\color{codegreen},
    keywordstyle=\color{magenta},
    numberstyle=\tiny\color{codegray},
    stringstyle=\color{codepurple},
    basicstyle=\ttfamily\footnotesize,
    breakatwhitespace=false,         
    breaklines=true,                 
    captionpos=b,                    
    keepspaces=true,                 
    numbers=left,                    
    numbersep=5pt,                  
    showspaces=false,                
    showstringspaces=false,
    showtabs=false,                  
    tabsize=2
}

\lstset{style=mystyle}

\title{JEE Project: Backend Development Report}
\author{Gemini CLI Assistant}
\date{May 21, 2026}

\begin{document}

\maketitle

\section{Introduction}
This report details the backend implementation of the JEE Student Management Project. The backend is built using Jakarta EE 10, leveraging JAX-RS for RESTful services, JPA for persistence, and EJB for business logic.

\section{Project Architecture}
The backend follows a layered architecture to ensure separation of concerns:
\begin{itemize}
    \item \textbf{Entity Layer}: JPA entities representing the database schema.
    \item \textbf{Repository Layer}: EJB-managed data access objects (DAOs).
    \item \textbf{Service/Security Layer}: Logic for password hashing and authentication.
    \item \textbf{Resource Layer}: JAX-RS endpoints providing the REST API.
\end{itemize}

\section{Functionality and Implementation}

\subsection{Authentication System}
\textbf{Functionality}: Secure user login and session token issuance.
\begin{itemize}
    \item \textbf{Backend Component}: \texttt{AuthResource.java} and \texttt{PasswordService.java}.
    \item \textbf{Function}: \texttt{login(User loginRequest)} in \texttt{AuthResource}.
    \item \textbf{Logic}:
    \begin{enumerate}
        \item Receives a JSON payload containing username and password.
        \item Calls \texttt{userRepository.findByUsername()} to retrieve the user.
        \item Uses \texttt{passwordService.verifyPassword()} (PBKDF2) to validate credentials.
        \item Returns a 200 OK response with a session token and user role if successful.
    \end{enumerate}
\end{itemize}

\subsection{Student Management (CRUD)}
\textbf{Functionality}: Full lifecycle management of student records.
\begin{itemize}
    \item \textbf{Backend Component}: \texttt{StudentResource.java} and \texttt{StudentRepository.java}.
    \item \textbf{Functions}:
    \begin{itemize}
        \item \textbf{Read}: \texttt{getAllStudents()} - Returns a list of all students using \texttt{studentRepository.findAll()}.
        \item \textbf{Create/Update}: \texttt{createStudent(Student student)} - Uses \texttt{studentRepository.save()} which performs either \texttt{persist()} or \texttt{merge()}.
        \item \textbf{Delete}: \texttt{deleteStudent(Long id)} - Uses \texttt{studentRepository.delete()} to remove a record by its primary key.
    \end{itemize}
\end{itemize}

\section{Data Persistence}
Persistence is managed via \texttt{persistence.xml} using the \texttt{CRUD\_PU} unit.
\begin{itemize}
    \item \textbf{Database}: H2 (In-memory) via WildFly's \texttt{ExampleDS}.
    \item \textbf{Schema Generation}: Set to \texttt{drop-and-create} to ensure a fresh environment on every deployment.
    \item \textbf{Initialization}: \texttt{import.sql} is automatically loaded to populate the database with default users and sample students.
\end{itemize}

\section{Security Implementation}
Passwords are never stored in plain text. We utilize the \texttt{Pbkdf2PasswordHash} mechanism from Jakarta Security.
\begin{itemize}
    \item \textbf{Algorithm}: PBKDF2WithHmacSHA256.
    \item \textbf{Iterations}: 2048.
    \item \textbf{Salt Size}: 32 bytes.
\end{itemize}

\end{document}
